\FigureLayoutDeclare{img/2010/guangdong_science/q18_fig1.png}{11.0cm}{24bp 20bp 41bp 30bp}

\chapter{2010年广东卷（理）}

\section{单选题}

\begin{problem}
若集合 \(A = \{x \mid -2 < x < 1\}\)，\(B = \{x \mid 0 < x < 2\}\)，则集合 \(A \cap B =\)（\quad）

\choices
    {\( \{x \mid -1 < x < 1\} \)}
    {\(\{x\mid -2 < x < 1\}\)}
    {\(\{x\mid -2 < x < 2\}\)}
    {\(\{x\mid 0 < x < 1\}\)}
\end{problem}

\begin{answer}
D
\end{answer}

\begin{solution}
集合交集中的元素必须同时满足两个集合的条件，因此 \(-2<x<1\) 且 \(0<x<2\)，从而得到 \(0<x<1\)，所以选 D.
\end{solution}

\begin{problem}
若复数 \(z_1 = 1 + \i\)，\(z_2 = 3 - \i\)，则 \( z_1 \cdot z_2 = \)（\quad）

\choices
    {\(4 + 2\i\)}
    {\(2 + \i\)}
    {\(2 + 2\i\)}
    {\(3 + \i\)}
\end{problem}

\begin{answer}
A
\end{answer}

\begin{solution}
计算得 \(z_1z_2=(1+\i)(3-\i)=3-\i+3\i-\i^2=4+2\i\)，所以选 A.
\end{solution}

\begin{problem}
若函数 \( f(x) = 3^x + 3^{-x} \) 与 \( g(x) = 3^x - 3^{-x} \) 的定义域均为 \( \R \)，则（\quad）

\choices
    {\(f(x)\) 与 \(g(x)\) 均为偶函数}
    {\(f(x)\) 为偶函数，\(g(x)\) 为奇函数}
    {\(f(x)\) 与 \(g(x)\) 均为奇函数}
    {\(f(x)\) 为奇函数，\(g(x)\) 为偶函数}
\end{problem}

\begin{answer}
B
\end{answer}

\begin{solution}
函数的定义域 \(\R\) 关于原点对称，且
\[
f(-x)=3^{-x}+3^x=f(x),\qquad g(-x)=3^{-x}-3^x=-g(x)
\]
因此 \(f(x)\) 是偶函数，\(g(x)\) 是奇函数，所以选 B.
\end{solution}

\begin{problem}
已知数列 \(\{a_{n}\}\) 为等比数列，\(S_{n}\) 是它的前 \(n\) 项和；若 \(a_{2} \cdot a_{3} = 2a_{1}\)，且 \(a_{4}\) 与 \(2a_{7}\) 的等差中项为 \(\frac{5}{4}\)，则 \(S_{5} =\)（\quad）

\choices
    {35}
    {33}
    {31}
    {29}
\end{problem}

\begin{answer}
C
\end{answer}

\begin{solution}
设等比数列的首项为 \(a_1\)，公比为 \(q\). 若 \(a_1=0\)，则所有项均为零，与等差中项为 \(\frac54\) 矛盾，故 \(a_1\ne0\). 由 \(a_2a_3=2a_1\) 得
\[
a_1^2q^3=2a_1,\qquad a_1q^3=2
\]
于是 \(a_4=a_1q^3=2\)，\(a_7=a_1q^6=(a_1q^3)q^3=2q^3\). 由 \(a_4\) 与 \(2a_7\) 的等差中项为 \(\frac54\)，得
\[
\frac{a_4+2a_7}{2}=\frac{2+4q^3}{2}=\frac54,\qquad q^3=\frac18,\qquad q=\frac12
\]
再由 \(a_1q^3=2\) 得 \(a_1=16\)，所以
\[
S_5=16\frac{1-(\frac12)^5}{1-\frac12}=31
\]
故选 C.
\end{solution}

\begin{problem}
“ \( m < \frac{1}{4} \) ”是一元二次方程 \( x^{2} + x + m = 0 \) 有实数解的（\quad）

\choices
    {充分非必要条件}
    {充分必要条件}
    {必要非充分条件}
    {非充分非必要条件}
\end{problem}

\begin{answer}
A
\end{answer}

\begin{solution}
一元二次方程有实数解的充要条件是判别式满足
\[
\Delta=1-4m\geqslant0,\qquad m\leqslant\frac14
\]
条件 \(m<\frac14\) 能推出 \(m\leqslant\frac14\)，但 \(m=\frac14\) 时方程仍有实数解，因此它是充分非必要条件，选 A.
\end{solution}

\begin{problem}
如图，\(\triangle ABC\) 为正三角形. \(AA' \parallel BB' \parallel CC'\)，\(CC' \perp\) 平面 \(ABC\)，且 \(3AA' = \frac{3}{2} BB' = CC' = AB\)，则多面体 \(ABC - A'B'C'\) 的正视图（也称主视图）是（\quad）

\choices
    {\choicebitmap[width=0.15\paperwidth]{img/2010/guangdong_science/q6_optA.png}}
    {\choicebitmap[width=0.15\paperwidth]{img/2010/guangdong_science/q6_optB.png}}
    {\choicebitmap[width=0.15\paperwidth]{img/2010/guangdong_science/q6_optC.png}}
    {\choicebitmap[width=0.15\paperwidth]{img/2010/guangdong_science/q6_optD.png}}
\end{problem}

\begin{answer}
D
\end{answer}

\begin{solution}
令 \(AB=1\). 由 \(AA'\parallel BB'\parallel CC'\) 且 \(CC'\perp\) 平面 \(ABC\)，可知三条线段均垂直于平面 \(ABC\). 由
\[
3AA'=\frac32BB'=CC'=AB=1
\]
得
\[
AA'=\frac13,\qquad BB'=\frac23,\qquad CC'=1
\]
按图中的正视方向投影，点 \(C\) 的投影位于 \(AB\) 的中点，故 \(CC'\) 在正视图中为中间的虚线；点 \(C'\) 的高度最高，且左端点 \(A'\) 的高度为 \(\frac13\)，低于右端点 \(B'\) 的高度 \(\frac23\). 于是正视图的外轮廓顶点依次为左端的 \(A'\)、中间最高的 \(C'\) 和右端的 \(B'\)，内部可见棱 \(A'B'\) 从左向右上升，符合选项 D. 因此选 D.
\end{solution}

\begin{problem}
已知随机变量 \(X\) 服从正态分布 \( N(3,1) \)，且 \( P(2 \leqslant X \leqslant 4) = 0.6826 \)，则 \( P(X > 4) = \)（\quad）

\choices
    {0.1588}
    {0.1587}
    {0.1586}
    {0.1585}
\end{problem}

\begin{answer}
B
\end{answer}

\begin{solution}
\(X\sim N(3,1)\) 表示均值为 \(3\)、标准差为 \(1\) 的正态分布，因此
\[
P(2\leqslant X\leqslant4)=P(3-1\leqslant X\leqslant3+1)=0.6826
\]
正态分布关于均值对称，所以两侧尾部概率相等，从而
\[
P(X>4)=\frac{1-0.6826}{2}=0.1587
\]
故选 B.
\end{solution}

\begin{problem}
为了迎接 2010 年广州亚运会，某大楼安装 \(5\text{ 个}\) 彩灯，它们闪亮的顺序不固定，每个彩灯闪亮只能是红，橙，黄，绿，蓝中的一种颜色，且这 \(5\text{ 个}\) 彩灯所闪亮的颜色各不相同. 记这 \(5\text{ 个}\) 彩灯有序地闪亮一次为一个闪烁，在每个闪烁中，每秒钟有且仅有一个彩灯闪亮，而相邻两个闪烁的时间间隔均为 \(5\text{ 秒}\). 如果要实现所有不同的闪烁，那么需要的时间至少是（\quad）

\choices
    {\(1205\text{ 秒}\)}
    {\(1200\text{ 秒}\)}
    {\(1195\text{ 秒}\)}
    {\(1190\text{ 秒}\)}
\end{problem}

\begin{answer}
C
\end{answer}

\begin{solution}
五个彩灯的颜色各不相同，五种颜色的闪亮顺序共有 \(5!=120\) 种，因此需要完成 \(120\) 个闪烁. 每个闪烁持续 \(5\) 秒，且相邻两个闪烁之间有 \(5\) 秒间隔，所以从第一个闪烁开始到最后一个闪烁结束所需的最短时间为
\[
120\times5+(120-1)\times5=1195\text{ 秒}
\]
故选 C.
\end{solution}

\section{填空题}

\begin{problem}
函数 \( f(x)=\lg(x-2) \) 的定义域是\fillinblank{}.
\end{problem}

\begin{answer}
\(x>2\)
\end{answer}

\begin{solution}
对数的真数必须为正，即 \(x-2>0\)，所以函数的定义域为 \(x>2\).
\end{solution}

\begin{problem}
若向量 \(\bs{a} = (1, 1, x)\)，\(\bs{b} = (1, 2, 1)\)，\(\bs{c} = (1, 1, 1)\)，满足条件 \( (\bs{c} - \bs{a}) \cdot (2\bs{b}) = -2 \)，则 \(x=\)\fillinblank{}.
\end{problem}

\begin{answer}
2
\end{answer}

\begin{solution}
由题意
\[
\bs{c}-\bs{a}=(0,0,1-x),\qquad 2\bs{b}=(2,4,2)
\]
因此
\[
(\bs{c}-\bs{a})\cdot(2\bs{b})=2(1-x)=-2,\qquad x=2
\]
\end{solution}

\begin{problem}
已知 \(a\)，\(b\)，\(c\) 分别是 \(\triangle ABC\) 的三个内角 \(A\)，\(B\)，\(C\) 所对的边. 若 \(a = 1\)，\(b = \sqrt{3}\)，\(A + C = 2B\)，则 \(\sin C = \)\fillinblank{}.
\end{problem}

\begin{answer}
1
\end{answer}

\begin{solution}
三角形内角和给出 \(A+B+C=\pi\)，结合 \(A+C=2B\) 得 \(3B=\pi\)，所以 \(B=\frac\pi3\). 由正弦定理
\[
\frac{a}{\sin A}=\frac{b}{\sin B}
\]
可得 \(\sin A=\frac{a\sin B}{b}=\frac12\). 又因为 \(A+C=\frac{2\pi}{3}\)，且 \(A,C>0\)，所以 \(0<A<\frac{2\pi}{3}\)，从而 \(A=\frac\pi6\)，\(C=\frac\pi2\). 因此 \(\sin C=1\).
\end{solution}

\begin{problem}
若圆心在 \(x\) 轴上，半径为 \(\sqrt{2}\) 的圆 \(O\) 位于 \(y\) 轴左侧，且与直线 \(x + y = 0\) 相切，则圆 \(O\) 的方程是\fillinblank{}.
\end{problem}

\begin{answer}
\((x+2)^2+y^2=2\)
\end{answer}

\begin{solution}
设圆心为 \(O(u,0)\). 由于圆位于 \(y\) 轴左侧，故 \(u<0\). 圆心到切线 \(x+y=0\) 的距离等于半径 \(\sqrt2\)，所以
\[
\frac{|u|}{\sqrt2}=\sqrt2,\qquad |u|=2
\]
结合 \(u<0\) 得 \(u=-2\)，故圆的方程为 \((x+2)^2+y^2=2\).
\end{solution}

\begin{problem}
某城市缺水问题比较突出，为了制定节水管理办法，对全市居民某年的月均用水量进行了抽样调查，其中 \(n\) 位居民的月均用水量分别为 \(x_{1}\)，\(\cdots\)，\(x_{n}\) （单位：\(\text{吨}\)）. 根据如图所示的程序框图，若 \(n = 2\)，且 \(x_{1}\)，\(x_{2}\) 分别为 \(1\)，\(2\)，则输出的结果 \(s\) 为\fillinblank{}.

\texfigure[max width=0.42\linewidth]{img_repaint/2010/guangdong_science/q13_fig1.tex}
\end{problem}

\begin{answer}
\(\frac{1}{4}\)
\end{answer}

\begin{solution}
程序开始时 \(s_1=0\)，\(s_2=0\)，\(i=1\). 当 \(i\leqslant n\) 时，程序先更新 \(s_1\) 和 \(s_2\)，再按当前的 \(i\) 计算 \(s\)，最后令 \(i=i+1\).

当 \(i=1\) 时，代入 \(x_1=1\)，得到 \(s_1=1\)，\(s_2=1\)，并计算
\[
s=\frac11\left(1-\frac11\times1^2\right)=0
\]
随后 \(i=2\). 当 \(i=2\) 时，代入 \(x_2=2\)，得到 \(s_1=3\)，\(s_2=5\)，并计算
\[
s=\frac12\left(5-\frac12\times3^2\right)=\frac14
\]
随后 \(i=3\)，此时 \(i\leqslant n\) 不成立，程序输出上一次计算得到的 \(s\)，所以 \(s=\frac14\).
\end{solution}

\section{选考题：填空题}

\begin{problem}
如图，\(AB\)，\(CD\) 是半径为 \(a\) 的圆 \(O\) 的两条弦，它们相交于 \(AB\) 的中点 \(P\)，\(PD = \frac{2a}{3}\)，\(\angle OAP = 30^\circ\)，则 \(CP = \)\fillinblank{}.
\end{problem}

\begin{answer}
\(\frac{9}{8}a\)
\end{answer}

\begin{solution}
因为弦 \(AB\) 的中点为 \(P\)，所以半径 \(OP\) 垂直于弦 \(AB\)，从而三角形 \(OAP\) 是直角三角形. 因此
\[
AP=OA\cos\angle OAP=a\cos30^\circ=\frac{\sqrt3}{2}a
\]
又因 \(P\) 是 \(AB\) 的中点，故 \(PB=AP=\frac{\sqrt3}{2}a\). 由相交弦定理
\[
PA\cdot PB=PC\cdot PD,\qquad PC=\frac{(\frac{\sqrt3}{2}a)^2}{\frac23a}=\frac98a
\]
\end{solution}

\begin{problem}
在极坐标系 \((\rho, \theta)\) （\(0 \leqslant \theta < 2\pi\)）中，曲线 \(\rho = 2\sin \theta\) 与 \(\rho \cos \theta = -1\) 的交点的极坐标为\fillinblank{}.
\end{problem}

\begin{answer}
\(\left(\sqrt{2},\frac{3\pi}{4}\right)\)
\end{answer}

\begin{solution}
将 \(\rho=2\sin\theta\) 代入 \(\rho\cos\theta=-1\)，得到
\[
2\sin\theta\cos\theta=\sin2\theta=-1
\]
在 \(0\leqslant\theta<2\pi\) 内，\(\sin2\theta=-1\) 给出 \(\theta=\frac{3\pi}{4}\) 或 \(\theta=\frac{7\pi}{4}\). 按极径定义取 \(\rho\geqslant0\)：当 \(\theta=\frac{3\pi}{4}\) 时，\(\rho=2\sin\frac{3\pi}{4}=\sqrt2\)；当 \(\theta=\frac{7\pi}{4}\) 时，\(\rho=-\sqrt2\)，不合要求. 所以交点的极坐标为 \(\left(\sqrt2,\frac{3\pi}{4}\right)\).
\end{solution}

\section{解答题}

\begin{problem}
已知函数 \( f(x) = A \sin (3x + \varphi) \) （\(A > 0, x \in (-\infty, +\infty), 0 < \varphi < \pi\)） 在 \( x = \frac{\pi}{12} \) 时取得最大值 4.

\begin{enumerate}
\item 求 \( f(x) \) 的最小正周期；
\item 求 \( f(x) \) 的解析式；
\item 若 \( f\left(\frac{2}{3}\alpha + \frac{\pi}{12}\right) = \frac{12}{5} \)，求 \( \sin \alpha \).
\end{enumerate}
\end{problem}

\begin{solution}
\begin{enumerate}
\item 函数 \(f(x)=A\sin(3x+\varphi)\) 的角频率为 \(3\)，所以最小正周期为
\[
T=\frac{2\pi}{3}
\]
\item 因为正弦函数的最大值为 \(1\)，且 \(A>0\)，由函数在 \(x=\frac{\pi}{12}\) 时取得最大值 \(4\) 可得 \(A=4\). 此时
\[
3\cdot\frac{\pi}{12}+\varphi=\frac\pi4+\varphi=\frac\pi2+2k\pi
\]
结合 \(0<\varphi<\pi\)，得到 \(\varphi=\frac\pi4\)，所以
\[
f(x)=4\sin\left(3x+\frac\pi4\right)
\]
\item 将题设自变量代入上式，得
\[
4\sin\left(3\left(\frac23\alpha+\frac\pi{12}\right)+\frac\pi4\right)=\frac{12}{5}
\]
化简得 \(4\sin\left(2\alpha+\frac\pi2\right)=\frac{12}{5}\)，即 \(\cos2\alpha=\frac35\). 因此
\[
\sin^2\alpha=\frac{1-\cos2\alpha}{2}=\frac15,\qquad \sin\alpha=\pm\frac{\sqrt5}{5}
\]
\end{enumerate}
\end{solution}

\begin{problem}
某食品厂为了检查一条自动包装流水线的生产情况. 随机抽取该流水线上的 \(40\text{ 件}\) 产品作为样本称出它们的重量（单位：\(\text{克}\)）. 重量的分组区间为 \((490, 495]\)，\((495, 500]\)，…，\((510, 515]\)，由此得到样本的频率分布直方图，如图所示.

\begin{enumerate}
\item 根据频率分布直方图，求重量超过 \(505\text{ 克}\) 的产品数量；
\item 在上述抽取的 \(40\text{ 件}\) 产品中任取 \(2\text{ 件}\)，设 \(Y\) 为重量超过 \(505\text{ 克}\) 的产品数量，求 \(Y\) 的分布列；
\item 从流水线上任取 \(5\text{ 件}\) 产品，求恰有 \(2\text{ 件}\) 产品合格的重量超过 \(505\text{ 克}\) 的概率.
\end{enumerate}

\texfigure[max width=0.50\linewidth]{img_repaint/2010/guangdong_science/q17_fig1.tex}
\end{problem}

\begin{solution}
\begin{enumerate}
\item 直方图中各组组距为 \(5\text{ 克}\)，所以重量超过 \(505\text{ 克}\) 的两组频率之和为
\[
(0.05+0.01)\times5=0.30
\]
因此重量超过 \(505\text{ 克}\) 的产品数量为 \(40\times0.30=12\text{ 件}\).
\item 从上述 \(40\) 件产品中任取 \(2\) 件，超过 \(505\text{ 克}\) 的有 \(12\) 件，其余有 \(28\) 件. 随机变量 \(Y\) 的可能取值为 \(0,1,2\)，且
\[
P(Y=0)=\frac{\mathrm C_{28}^{2}}{\mathrm C_{40}^{2}},\qquad P(Y=1)=\frac{\mathrm C_{28}^{1}\mathrm C_{12}^{1}}{\mathrm C_{40}^{2}},\qquad P(Y=2)=\frac{\mathrm C_{12}^{2}}{\mathrm C_{40}^{2}}
\]
其分布列为

\begin{center}
\begin{tabular}{c|ccc}
\(Y\) & \(0\) & \(1\) & \(2\) \\
\hline
\(P\) & \(\dfrac{\mathrm C_{28}^{2}}{\mathrm C_{40}^{2}}\) & \(\dfrac{\mathrm C_{28}^{1}\mathrm C_{12}^{1}}{\mathrm C_{40}^{2}}\) & \(\dfrac{\mathrm C_{12}^{2}}{\mathrm C_{40}^{2}}\)
\end{tabular}
\end{center}
\item 由样本频率估计流水线上任取一件产品的重量超过 \(505\text{ 克}\) 的概率为 \(p=\frac{12}{40}=\frac3{10}\)，不超过 \(505\text{ 克}\) 的概率为 \(1-p=\frac7{10}\). 任取 \(5\) 件时，恰有 \(2\) 件超过 \(505\text{ 克}\) 的概率为
\[
\mathrm C_{5}^{2}\left(\frac3{10}\right)^2\left(\frac7{10}\right)^3=\frac{3087}{10000}
\]
\end{enumerate}
\end{solution}

\begin{problem}
如图，\(\widehat{AEC}\) 是半径为 \(a\) 的半圆，\(AC\) 为直径，点 \(E\) 为 \(\widehat{AC}\) 的中点，点 \(B\) 和点 \(C\) 为线段 \(AD\) 的三等分点，平面 \(AEC\) 外一点 \(F\) 满足 \(FB = FD = \sqrt{5}a\)，\(EF = \sqrt{6}a\).

\begin{enumerate}
\item 证明： \(EB \perp FD\)；
\item 已知点 \(Q\)，\(R\) 分别为线段 \(FE\)，\(FB\) 上的点，使得 \(FQ = \frac{2}{3} FE\)，\(FR = \frac{2}{3} FB\)，求平面 \(BED\) 与平面 \(RQD\) 所成二面角的正弦值.
\end{enumerate}

\bitmapfigure[max width=0.36\linewidth]{img/2010/guangdong_science/q18_fig1.png}
\end{problem}

\begin{solution}
\begin{enumerate}
\item 取 \(A=(-a,0,0)\)，\(C=(a,0,0)\)，\(E=(0,a,0)\)，则 \(B=(0,0,0)\)，\(D=(2a,0,0)\). 设 \(F=(u,v,w)\). 由 \(FB=FD\) 得
\[
u^2+v^2+w^2=(u-2a)^2+v^2+w^2,\qquad u=a
\]
再由 \(FB=\sqrt5a\) 和 \(FE=\sqrt6a\)，分别得
\[
v^2+w^2=4a^2,\qquad a^2+(v-a)^2+w^2=6a^2
\]
将第一式代入第二式，得 \(a^2+4a^2-2av+a^2=6a^2\)，所以 \(v=0\)，进而 \(w^2=4a^2\). 令 \(\varepsilon\in\{1,-1\}\)，则可写作 \(F=(a,0,2\varepsilon a)\). 因此 \(\overrightarrow{EB}=(0,-a,0)\)，\(\overrightarrow{FD}=(a,0,-2\varepsilon a)\)，从而
\[
\overrightarrow{EB}\cdot\overrightarrow{FD}=0
\]
所以 \(EB\perp FD\).
\item 由线段内分关系
\[
Q=F+\frac23(E-F)=\left(\frac a3,\frac{2a}{3},\frac{2\varepsilon a}{3}\right),\qquad R=F+\frac23(B-F)=\left(\frac a3,0,\frac{2\varepsilon a}{3}\right)
\]
平面 \(BED\) 为 \(z=0\)，可取其法向量 \(\bs n_1=(0,0,1)\). 又
\[
\overrightarrow{RQ}=(0,\tfrac{2a}{3},0),\qquad \overrightarrow{RD}=(\tfrac{5a}{3},0,-\tfrac{2\varepsilon a}{3}),
\]
所以平面 \(RQD\) 的法向量可取
\[
\bs n_2=\overrightarrow{RQ}\times\overrightarrow{RD}\sim(2\varepsilon,0,5)
\]
设两平面所成的锐二面角为 \(\theta\)，则
\[
\sin\theta=\frac{|\bs n_1\times\bs n_2|}{|\bs n_1||\bs n_2|}=\frac2{\sqrt{29}}=\frac{2\sqrt{29}}{29}
\]
\end{enumerate}
\end{solution}

\begin{problem}
某营养师要为某个儿童预定午餐和晚餐. 已知 \(1\text{ 个单位}\) 的午餐含 \(12\text{ 个单位}\) 的碳水化合物，\(6\text{ 个单位}\) 的蛋白质和 \(6\text{ 个单位}\) 的维生素 \(\mathrm C\)；\(1\text{ 个单位}\) 的晚餐含 \(8\text{ 个单位}\) 的碳水化合物，\(6\text{ 个单位}\) 的蛋白质和 \(10\text{ 个单位}\) 的维生素 \(\mathrm C\). 另外，该儿童这两餐需要的营养中至少含 \(64\text{ 个单位}\) 的碳水化合物，\(42\text{ 个单位}\) 的蛋白质和 \(54\text{ 个单位}\) 的维生素 \(\mathrm C\). 如果 \(1\text{ 个单位}\) 的午餐，晚餐的费用分别是 \(2.5\text{ 元}\) 和 \(4\text{ 元}\)，那么要满足上述的营养要求，并且花费最少，应当为该儿童分别预订多少个单位的午餐和晚餐？
\end{problem}

\begin{solution}
设午餐和晚餐分别预订 \(x\)，\(y\) 个单位，则 \(x\geqslant0\)，\(y\geqslant0\)，并满足
\[
3x+2y\geqslant16,\qquad x+y\geqslant7,\qquad 3x+5y\geqslant27
\]
费用为 \(C=2.5x+4y\). 将费用改写为
\[
C=\frac14(x+y)+\frac34(3x+5y)
\]
由营养约束可得
\[
C\geqslant\frac14\times7+\frac34\times27=22
\]
当 \(x=4\)，\(y=3\) 时，三项营养量分别为 \(72\)，\(42\)，\(54\)，满足全部要求，且费用为 \(2.5\times4+4\times3=22\text{ 元}\). 因此应预订 \(4\) 个单位的午餐和 \(3\) 个单位的晚餐，最低费用为 \(22\text{ 元}\).
\end{solution}

\begin{problem}
已知双曲线 \(\frac{x^2}{2} - y^2 = 1\) 的左、右顶点分别为 \(A_1\)，\(A_2\)，点 \(P(x_1, y_1)\)，\(Q(x_1, -y_1)\) 是双曲线上不同的两个动点.

\begin{enumerate}
\item 求直线 \(A_{1}P\) 与 \(A_{2}Q\) 交点的轨迹 \(E\) 的方程；
\item 若过点 \(H(0, h)\)（\(h > 1\)）的两条直线 \(l_{1}\) 和 \(l_{2}\) 与轨迹 \(E\) 都只有一个交点，且 \(l_{1} \perp l_{2}\)，求 \(h\) 的值.
\end{enumerate}
\end{problem}

\begin{solution}
\begin{enumerate}
\item 设交点为 \(M(x,y)\). 直线 \(A_1P\) 和 \(A_2Q\) 的方程分别为
\[
y=\frac{y_1}{x_1+\sqrt2}(x+\sqrt2),\qquad y=-\frac{y_1}{x_1-\sqrt2}(x-\sqrt2)
\]
因 \(P\)，\(Q\) 是不同的点，故 \(y_1\ne0\). 联立两式，得
\[
\frac{y_1(x+\sqrt2)}{x_1+\sqrt2}=-\frac{y_1(x-\sqrt2)}{x_1-\sqrt2}
\Longrightarrow (x+\sqrt2)(x_1-\sqrt2)=-(x-\sqrt2)(x_1+\sqrt2)
\Longrightarrow xx_1=2
\]
所以 \(x=\frac2{x_1}\). 代回第一式得 \(y=\frac{\sqrt2y_1}{x_1}\). 由 \(P\) 在双曲线上，\(\frac{x_1^2}{2}-y_1^2=1\)，且 \(y_1\ne0\)，可知 \(\lvert x_1\rvert>\sqrt2\)，所以
\[
\frac{x^2}{2}+y^2=\frac2{x_1^2}+\frac{2y_1^2}{x_1^2}=\frac{2(1+y_1^2)}{x_1^2}=1
\]
因此轨迹 \(E\) 的方程为 \(\frac{x^2}{2}+y^2=1\)，且 \(0<|x|<\sqrt2\).
\item 过 \(H(0,h)\) 的直线可设为 \(y=mx+h\). 代入轨迹方程，得到关于 \(x\) 的二次方程
\[
\left(m^2+\frac12\right)x^2+2mhx+h^2-1=0
\]
由第（1）问，轨迹 \(E\) 是椭圆 \(\frac{x^2}{2}+y^2=1\) 去掉 \((0,1)\)，\((0,-1)\)，\((\sqrt2,0)\)，\((-\sqrt2,0)\) 后的曲线. 因此，直线与 \(E\) 只有一个交点时，既可能与椭圆相切，也可能经过一个被去掉的端点而与椭圆相交于另一个仍属于 \(E\) 的点.

先考虑相切情形. 判别式为零，从而
\[
(2mh)^2-4\left(m^2+\frac12\right)(h^2-1)=0,\qquad m^2=\frac{h^2-1}{2}
\]
因 \(h>1\)，此时 \(m\ne0\)，而椭圆在 \((0,\pm1)\) 处的切线为水平线、在 \((\pm\sqrt2,0)\) 处的切线为竖直线，所以切点属于 \(E\).
两条切线的斜率为相反数，若两条直线均为切线且互相垂直，则
\[
-m^2=-1,\qquad h^2=3,\qquad h=\sqrt3.
\]

再考虑相交情形. 直线 \(x=0\) 虽经过被去掉的点 \((0,1)\) 和 \((0,-1)\)，但它与 \(E\) 没有交点；因此，能使相交的两个椭圆交点中恰有一个属于 \(E\) 的非竖直直线，只能经过被去掉的两个顶点 \((\sqrt2,0)\) 和 \((-\sqrt2,0)\). 过 \(H\) 与这两个点分别作直线，其斜率为
\[
\frac{h}{\sqrt2},\qquad -\frac{h}{\sqrt2}.
\]
这两条直线均经过一个不属于 \(E\) 的顶点，且与椭圆的另一个交点属于 \(E\)，故各自与 \(E\) 只有一个交点. 若它们互相垂直，则
\[
-\frac{h^2}{2}=-1,\qquad h=\sqrt2.
\]

最后考虑一条切线和一条上述相交直线的情形. 两条直线互相垂直时斜率符号相反，故
\[
\sqrt{\frac{h^2-1}{2}}\cdot\frac{h}{\sqrt2}=1,\qquad h^2(h^2-1)=4.
\]
令 \(u=h^2\)，则 \(u>1\) 且 \(u^2-u-4=0\)，所以
\[
h=\sqrt{\frac{1+\sqrt{17}}{2}}.
\]
综上，\(h\) 的值为
\[
\boxed{\sqrt2},\qquad \boxed{\sqrt3},\qquad
\boxed{\sqrt{\frac{1+\sqrt{17}}{2}}}.
\]
\end{enumerate}
\end{solution}

\begin{problem}
设 \(A(x_{1},y_{1})\)，\(B(x_{2},y_{2})\) 是平面直角坐标系 \(xOy\) 上的两点. 现定义由点 \(A\) 到点 \(B\) 的一种折线距离 \(\rho(A,B)\) 为 \(\rho(A,B) = |x_2 - x_1| + |y_2 - y_1|\). 对于平面 \(xOy\) 上给定的不同两点 \(A(x_{1},y_{1})\)，\(B(x_{2},y_{2})\).

\begin{enumerate}
\item 若点 \(C(x, y)\) 是平面 \(xOy\) 上的点，试证明 \(\rho(A,C) + \rho(C,B) \geqslant \rho(A,B)\)；
\item 在平面 \(xOy\) 上是否存在点 \(C(x, y)\)，同时满足
\begin{circlelist}
\item \(\rho(A,C) + \rho(C,B) = \rho(A,B)\)；
\item \(\rho(A,C) = \rho(C,B)\).
\end{circlelist}
若存在，请求出所有符合条件的点；若不存在，请予以证明.
\end{enumerate}
\end{problem}

\begin{solution}
\begin{enumerate}
\item 由绝对值三角不等式，分别有
\[
|x-x_1|+|x_2-x|\geqslant|x_2-x_1|,\qquad |y-y_1|+|y_2-y|\geqslant|y_2-y_1|
\]
两式相加即得 \(\rho(A,C)+\rho(C,B)\geqslant\rho(A,B)\).
\item 第一问中的等号成立，当且仅当 \(x\) 位于 \(x_1\)，\(x_2\) 之间且 \(y\) 位于 \(y_1\)，\(y_2\) 之间，即
\[
\min(x_1,x_2)\leqslant x\leqslant\max(x_1,x_2),\qquad \min(y_1,y_2)\leqslant y\leqslant\max(y_1,y_2)
\]
在此条件下，令 \(u=|x_2-x_1|\)，\(v=|y_2-y_1|\)，则
\[
\rho(A,C)=|x-x_1|+|y-y_1|,\qquad \rho(C,B)=u+v-\rho(A,C)
\]
故第二个等式等价于 \(|x-x_1|+|y-y_1|=\frac{u+v}{2}\). 因此所有符合条件的点组成矩形
\[
[\min(x_1,x_2),\max(x_1,x_2)]\times[\min(y_1,y_2),\max(y_1,y_2)]
\]
内与直线 \(|x-x_1|+|y-y_1|=\frac{u+v}{2}\) 的交集. 具体地，令
\[
t\in\left[\max\left(0,\frac{u-v}{2}\right),\min\left(u,\frac{u+v}{2}\right)\right]
\]
并记 \(\varepsilon_x=\operatorname{sgn}(x_2-x_1)\)，\(\varepsilon_y=\operatorname{sgn}(y_2-y_1)\)，则所有解为
\[
C=\left(x_1+\varepsilon_x t,\ y_1+\varepsilon_y\left(\frac{u+v}{2}-t\right)\right)
\]
当 \(x_1\ne x_2\) 且 \(y_1\ne y_2\) 时，这是一条线段；当两点有一个坐标相同时，线段退化为唯一的中点 \(C\left(\frac{x_1+x_2}{2},\frac{y_1+y_2}{2}\right)\).
\end{enumerate}
\end{solution}
